% TOP_PROBLEM: {"id": 2302011, "title": "Research Problems in Function Theory — Problem 2.11", "category_id": 8, "category": {"id": 8, "name": "analysis", "display_name": "Analysis", "description": "Limits, continuity, calculus, and function theory.", "slug": "analysis", "order_index": 8, "created_at": "2026-07-31T15:26:25.671Z"}, "status": "open", "problem_number": "AMR-022-2011", "set_id": 13}
% FIRST_POSTED: 2026-10-02
% ENTRY_KIND: statement_only
% UnsolvedMath ID 2302011; source catalogue code AMR-022-2011
% !TeX program = lualatex
% Definition and sources only; this file is not an LLM solution attempt.
\documentclass[11pt,a4paper]{article}
\usepackage[margin=25mm]{geometry}
\usepackage{fontspec}
\setmainfont{lmroman10-regular.otf}[BoldFont=lmroman10-bold.otf,
 ItalicFont=lmroman10-italic.otf,BoldItalicFont=lmroman10-bolditalic.otf]
\usepackage{amsmath,amssymb,mathtools}
\usepackage{unicode-math}
\setmathfont{latinmodern-math.otf}
\AtBeginDocument{\let\setminus\smallsetminus}
\usepackage{xurl,hyperref}
\hypersetup{colorlinks=true,urlcolor=blue}
\setcounter{secnumdepth}{0}
\setlength{\parindent}{0pt}
\setlength{\parskip}{5pt}
\setlength{\emergencystretch}{3em}
\begin{document}
\section*{Research Problems in Function Theory — Problem 2.11}\label{2302011}
{\small UnsolvedMath ID 2302011 · AMR-022-2011 · Analysis · AMR Open Problem Lists}
\subsection{Definitions and mathematical statement}
Let $0<\lambda_1<\lambda_2<\cdots$ be integers and let
\[
f(z)=a_0+\sum_{n\ge1}a_nz^{\lambda_n},\qquad a_n\ne0,
\qquad \sum_{n\ge1}\frac1{\lambda_n}<\infty,
\]
be a transcendental entire function. Thus the Taylor expansion has Fej\'er gaps; convergence is required on the whole complex plane, not just in the unit disc. Set
$M(r,f)=\max_{|z|=r}|f(z)|$ and $m_0(r,f)=\min_{|z|=r}|f(z)|$.
The problem has two assertions. First, can $f$ have a finite asymptotic value: a complex number $a$ and a continuous curve $\gamma:[0,\infty)\to\mathbb C$ such that $|\gamma(t)|\to\infty$ and $f(\gamma(t))\to a$? The proposed answer is no. Second, must
\[
\limsup_{r\to\infty}\frac{\log m_0(r,f)}{\log M(r,f)}=1?
\]
Here $\log m_0=-\infty$ at radii containing a zero; the denominator is positive for sufficiently large $r$. The reverse inequality is automatic from $m_0\le M$. The source records results under stronger gap assumptions; the summability assumption alone is the intended hypothesis in both questions.
\subsection{Short English statement}
Do summable reciprocal Taylor exponents exclude every finite asymptotic value and force near equality of minimum and maximum modulus along some large circles?
\subsection{Sources}
\begin{itemize}
\item[S1] ULAM.AI and the UnsolvedMath Contributors, supplied problems.json, record 2302011 (AMR-022-2011), AMR Open Problem Lists. Catalogue identity and source transcription; CC BY 4.0.
\url{https://huggingface.co/datasets/ulamai/UnsolvedMath}.
\item[S2] Hayman - Research Problems in Function Theory (2018), Problem 2.11. Original source indexed by AMR-022-2011.
\url{https://arxiv.org/abs/1809.07200}.
\end{itemize}
\end{document}
